\section{The measurement chain: how the pieces fit together}
\label{sec:chain}

Before the formula-by-formula derivations, this section tells the story
once, end to end, so that each later section can be read as a close-up of
one link in a chain you have already seen whole.

\subsection{From a magnitude to a spectrum}

The user gives the program one brightness number --- ``the target is
magnitude 12.3 in $V$'' --- but a detector does not respond to
magnitudes; it responds to photons per wavelength. The first task is
therefore to turn one number into a full spectral energy distribution.
This is what the \emph{template} is for: an absolutely calibrated
spectrum of a star of similar type supplies the \emph{shape}
$F_\lambda(\lambda)$, and the magnitude supplies the \emph{level}. The
connection between the two is the synthetic-photometry machinery of
Sect.~\ref{sec:radiometry}: the program computes what magnitude the
template \emph{would} have in the user's reference band, and rescales the
whole spectrum by the ratio to the requested magnitude. Because the
rescaling is total, the template's own flux units are immaterial --- only
its shape and its internal calibration matter --- which is why CALSPEC
standards, surface-brightness planet atlases and arbitrary-unit
supernova models can all serve as templates
(Sect.~\ref{sec:radiometry}). The one number that anchors this step is
the catalogue's $m_{v0}$, the magnitude the template file itself
represents; it is the fulcrum on which the whole calibration turns.

\subsection{From a spectrum above the atmosphere to electrons in a pixel}

The calibrated spectrum then runs the gauntlet of the light path, each
element a multiplicative wavelength-dependent factor inside one integral,
Eq.~(\ref{eq:rate}):

\begin{center}
$F_\lambda$ (target, top of atmosphere)
$\;\rightarrow\; T_{\rm atm}(\lambda)^{X}$ (Sect.~\ref{sec:atmosphere})
$\;\rightarrow\; A\,\eta\,O_\lambda$ (telescope and optics)
$\;\rightarrow\; T_\lambda$ (filter)
$\;\rightarrow\; Q_\lambda$ (detector QE)
$\;\rightarrow\;$ electrons s$^{-1}$
\end{center}

The atmosphere enters twice more, in subtler roles than mere absorption:
it \emph{blurs} (the seeing PSF of Sect.~\ref{sec:psf}, which decides how
much light a finite aperture or slit captures and how bright the peak
pixel gets), it \emph{disperses} (differential refraction,
Sect.~\ref{sec:dispersion}, which drains blue light out of a misaligned
slit), and it \emph{twinkles} (scintillation,
Sect.~\ref{sec:scintillation}, a noise proportional to the signal
itself). And the sky above the site is not dark: the sky model of
Sect.~\ref{sec:sky} --- airglow, zodiacal light, starlight, moonlight,
twilight, or simply the user's SQM reading --- adds its own electrons to
every aperture and every resolution element, through the same instrument
chain.

\subsection{From electrons to a decision}

The last link converts electron counts into the quantities an observer
decides with. The CCD equation, Eq.~(\ref{eq:snr}), combines the source
against every accumulated noise: its own shot noise, the sky's, dark
current, read noise per pixel, scintillation, quantization. Solved
forward it yields the S/N of a given exposure; inverted
(Sect.~\ref{sec:stack}) it yields the exposure for a requested S/N, and,
when that exposure would saturate the brightest pixel --- which the
program predicts independently, from the unextracted PSF peak --- the
stack planner divides it into feasible frames. Spectroscopy adds one
more translation: from S/N per resolution element to the
equivalent-width uncertainty $\sigma(\mathrm{EW})$ of
Sect.~\ref{sec:ew}, which states directly which spectral lines a night
can measure. Everything downstream of the electron counts is exact
bookkeeping; everything upstream is physics with stated approximations
--- the map of which is drawn in Sect.~\ref{sec:expectations}.

\subsection{Where each ingredient comes from}

Every factor in the chain is data the user controls, and each has a
documented file format (User Guide, Sect.~6): the template from the
catalogue; the filter transmission and zero point from an SVO profile
(downloaded or synthesised from a measured curve,
Sect.~\ref{sec:filters}); the QE, optics response and gain table from
the camera's datasheet; the atmospheric transmission from a site FITS
table or the built-in broad-band fallback; the sky level from the
physical model or a hand-held SQM meter. The design intent is that
\emph{nothing physical is hidden in the code}: what the program assumes,
a file or a visible control states, and the run log in the Results
window records.
